\section{The primitive-sixth-root atom}
\label{sec:one-step}

We use the geometric atomic \(F\)-bundle invariant of
Katzarkov--Kontsevich--Pantev--Yu \cite{KKPYY}.  Only three features enter:
its values lie in the free abelian group on atom-equivalence classes, it is
additive under blowups, and a rank-\(r\) projective bundle contributes \(r\)
copies of the base.

\begin{definition}
For a geometric atom \(\alpha\), let \(\FM(\alpha)\) be the multiset of
formal-monodromy eigenvalues on its Levelt--Turrittin regular-singular
factors, after the exponential parts have been separated.  Define
\[
 \nu_6(\alpha)=
 \operatorname{mult}_{e^{\pi i/3}}\FM(\alpha)
 +\operatorname{mult}_{e^{-\pi i/3}}\FM(\alpha).
\]
Extend \(\nu_6\) additively to the free atom group.
\end{definition}

The definition uses eigenvalues in \(\C^\times\), not chosen logarithms.
Integral powers of the loop coordinate shift residue exponents by integers
and leave \(\nu_6\) unchanged.

\begin{lemma}[Formal isomonodromy]
\label{lem:formal-isomonodromy}
Along a connected component of the reduced unramified spectral cover, the
formal-monodromy eigenvalue multiset of an atom is constant.
\end{lemma}

\begin{proof}
After finite formal ramification and separation of exponential factors, one
block has the form
\[
 \nabla=d-d\Phi(s,u)+R(s)\frac{du}{u}
       +\sum_i\Gamma_i(s,u)\,ds_i,
\]
where \(\Gamma_i\) is formal-regular after the exponential term is removed.
The coefficient of \(du/u\wedge ds_i\) in \(\nabla^2=0\) is
\[
 \partial_{s_i}R=[\Gamma_i(s,0),R].
\]
Hence the characteristic polynomial of \(R\), and therefore that of
\(\exp(2\pi iR)\), is locally constant.  Descent only permutes the separated
blocks.
\end{proof}

\begin{proposition}[Cubic and low-dimensional atoms]
\label{prop:atom-values}
Let \(X\) be a smooth cubic threefold.
\begin{enumerate}[label=\textup{(\roman*)}]
\item The cubic atom \(\alpha_X\) satisfies \(\nu_6(\alpha_X)>0\).
\item Every atom represented by a point, a curve, or a smooth projective
surface has \(\nu_6=0\).
\end{enumerate}
\end{proposition}

\begin{proof}
Cai's Proposition~6 gives formal residue classes \(\pm1/6\) in the
rank-two zero-eigenvalue block of the big cubic quantum connection
\cite{Cai}.  The same block is the unsplit cubic zero atom in the maximal
big-quantum \(F\)-bundle of \cite[Example~6.21]{KKPYY}; scalar extension does
not change its formal monodromy.  This proves (i).

For a minimal surface with nef canonical class, the parity-corrected
connection has nilpotent residue by \cite[Claim~6.15]{KKPYY}.  Undoing the
parity correction permits only eigenvalues \(1\) and \(-1\).  Ruled surfaces
reduce by the projective-bundle formula to curves, and nonminimal surfaces
reduce by point blowups to their minimal models.  Points and curves likewise
have only integral or half-integral residues.  The classification of smooth
projective surfaces exhausts the possibilities and proves (ii).
\end{proof}

\begin{proof}[Proof of Theorem~\ref{thm:every-cubic}]
Let \(Y=X\times\PP^1\).  The projective-bundle formula gives
\[
 \CFF(Y)=2\CFF(X),
\]
so Proposition~\ref{prop:atom-values} gives \(\nu_6(Y)>0\).

Suppose that \(Y\) were rational.  Weak factorization of a birational map
\(Y\dashrightarrow\PP^4\) gives, in the free abelian atom group,
\[
 \CFF(Y)+\sum_i(r_i-1)\CFF(Z_i)
 =\CFF(\PP^4)+\sum_j(s_j-1)\CFF(W_j),
\]
where every smooth center \(Z_i,W_j\) has dimension at most two.  Apply
\(\nu_6\).  Proposition~\ref{prop:atom-values}(ii) kills every center term,
and the projective-bundle formula over a point gives
\(\nu_6(\PP^4)=0\).  The left-hand side retains \(\nu_6(Y)>0\), a
contradiction.
\end{proof}

The proof uses the free atom group, not a signed informal block ledger.  In
particular the two copies supplied by \(\PP^1\) are positive multiplicities
and cannot cancel one another.
